\section{Data Format and Encoding Strategies}
\label{sec:data}

\subsection{Data Format}

Following the original DeepChrInteract~\cite{deepchrinteract_doc}, we use labelled
chromatin interaction data derived from Hi-C experiments.
The dataset for each cell type consists of four plain-text files:

\begin{itemize}[noitemsep]
  \item \texttt{seq.anchor1.pos.txt} --- enhancer sequences from interacting loci (positive);
  \item \texttt{seq.anchor2.pos.txt} --- corresponding promoter sequences (positive);
  \item \texttt{seq.anchor1.neg2.txt} --- enhancer sequences from non-interacting loci (negative);
  \item \texttt{seq.anchor2.neg2.txt} --- corresponding promoter sequences (negative).
\end{itemize}

Each line contains one DNA sequence of fixed length $L$ (default $L = 10{,}001$~bp)
over the alphabet $\{A, C, G, T, N\}$, where $N$ denotes an unresolved nucleotide.
Positive and negative samples are balanced at dataset construction time following
standard practice~\cite{speid}.
The dataset is split into training (80\%), validation (10\%), and test (10\%) sets
by stratified random sampling with a fixed seed.

\subsection{Encoding Strategy A: One-Hot}

One-hot encoding represents each nucleotide as a binary indicator vector of
dimension five (including $N$):
\begin{equation}
  \phi_{\mathrm{oh}}(c) =
  \begin{cases}
    \mathbf{e}_1 & c = A \\
    \mathbf{e}_2 & c = C \\
    \mathbf{e}_3 & c = G \\
    \mathbf{e}_4 & c = T \\
    \mathbf{e}_5 & c = N
  \end{cases}
  \in \{0,1\}^5,
\end{equation}
where $\mathbf{e}_k$ is the $k$-th standard basis vector.
A sequence of length $L$ is encoded as a matrix $\mathbf{X} \in \{0,1\}^{5 \times L}$,
which is treated as a 5-channel 1D signal and processed by 1D convolution layers.
The encoding is computed on-the-fly from the raw text files via a vectorised
NumPy kernel during data loading, with no disk-cached intermediate representation.

\paragraph{Design note.}
The original implementation converted one-hot matrices to $[0,1]$ float PNG images
(via \texttt{imageio}) and then read them back with Keras's \texttt{ImageDataGenerator}.
This incurred two conversion round-trips (float $\to$ uint8 $\to$ float), introduced
8-bit quantisation error, generated tens of gigabytes of temporary files, and created
a tight coupling between the data pipeline and the image loading subsystem.
DeepChrInteract-v2 eliminates this entirely: one-hot tensors are constructed in
\texttt{\_\_getitem\_\_} and never written to disk.

\subsection{Encoding Strategy B: \texorpdfstring{$k$}{k}-mer Tokenisation}

To capture local sequence context beyond individual nucleotides, we tile a sliding window
of width $k=6$ over each sequence, yielding $L - k + 1$ overlapping $k$-mers.
Since there are $4^6 = 4{,}096$ valid $k$-mers over $\{A,C,G,T\}$, plus a special
\texttt{null} token for windows containing $N$, the vocabulary size is 4,097.
Each $k$-mer is mapped to an integer index, and the integer sequence is fed to a
trainable \texttt{nn.Embedding} layer of output dimension $d_e = 128$.
This replicates the original tokenisation scheme, extended with the DNA2Vec
pre-trained weight matrix~\cite{deepchrinteract_doc} as an optional initialisation.

\subsection{Encoding Strategy C: Pre-Trained DNA LLM}

Pre-trained DNA foundation models provide contextualised nucleotide representations
that encode co-evolutionary constraints and regulatory signals learned from
large genomic corpora.
We integrate four models as optional encoders:

\begin{description}[noitemsep]
  \item[DNABERT~\cite{dnabert}]
    A 12-layer BERT model pre-trained on the human reference genome
    with 6-mer tokenisation. We use the publicly available weights from
    \texttt{DNABERT-6} and either freeze all parameters (as a fixed feature extractor)
    or fine-tune the top $K$ transformer layers together with the EPI classification head.

  \item[DNABERT-2~\cite{dnabert2}]
    An improved foundation model using BPE tokenisation and a Flash-Attention backbone,
    trained on 135 species. BPE naturally handles the variable-length motif structure
    of regulatory sequences and reduces sequence length relative to $k$-mer tokenisation.

  \item[Nucleotide Transformer~\cite{nucleotide_transformer}]
    A family of models (50M--2.5B parameters) pre-trained on 3,202 human genomes.
    We adopt the 500M multi-species variant for its balance of capacity and GPU memory
    footprint.

  \item[HyenaDNA~\cite{hyenadna}]
    A sub-quadratic architecture based on Hyena operators that processes up to
    $1 \times 10^6$ tokens at single-nucleotide resolution.
    Given the sequence length of $L = 10{,}001$~bp per region in our dataset,
    HyenaDNA can encode the full window without downsampling, making it the
    most natural fit for our task.
\end{description}

For all LLM encoders, the full sequence is passed through the model and the
\texttt{[CLS]} token representation (or the mean of all token representations
for models without a \texttt{[CLS]} token) is used as the fixed-length
representation $\mathbf{z} \in \mathbb{R}^{d}$.
This vector is subsequently passed to the fusion and classification modules
described in \Cref{sec:methods}.

\subsection{Data Pipeline Implementation}

\Cref{alg:dataset} summarises the data loading pipeline.
All encoding strategies share a common \texttt{EPIDataset} interface that returns
a tuple $(\mathbf{x}_e, \mathbf{x}_p, y)$ of enhancer encoding, promoter encoding,
and binary label, compatible with the standard PyTorch \texttt{DataLoader}.

\begin{algorithm}[H]
\caption{EPIDataset \texttt{\_\_getitem\_\_}}\label{alg:dataset}
\begin{algorithmic}[1]
\Require Index $i$, encoding mode $m \in \{\text{onehot}, \text{kmer}, \text{llm}\}$
\State Read raw sequences: $s_e \leftarrow \texttt{enhancer}[i]$,
       $s_p \leftarrow \texttt{promoter}[i]$, label $y \leftarrow \texttt{label}[i]$
\If{$m = \text{onehot}$}
  \State $\mathbf{x}_e \leftarrow \phi_{\mathrm{oh}}(s_e) \in \mathbb{R}^{5 \times L}$
  \State $\mathbf{x}_p \leftarrow \phi_{\mathrm{oh}}(s_p) \in \mathbb{R}^{5 \times L}$
\ElsIf{$m = \text{kmer}$}
  \State $\mathbf{x}_e \leftarrow \texttt{tokenize}(s_e) \in \mathbb{Z}^{L-5}$
  \State $\mathbf{x}_p \leftarrow \texttt{tokenize}(s_p) \in \mathbb{Z}^{L-5}$
\ElsIf{$m = \text{llm}$}
  \State $\mathbf{x}_e \leftarrow \texttt{LLM.encode}(s_e) \in \mathbb{R}^{d}$ \quad\textit{(pre-computed)}
  \State $\mathbf{x}_p \leftarrow \texttt{LLM.encode}(s_p) \in \mathbb{R}^{d}$
\EndIf
\State \Return $(\mathbf{x}_e,\; \mathbf{x}_p,\; y)$
\end{algorithmic}
\end{algorithm}
